\chapter{Open Problems and the Remaining Physical Burden}\label{ch:open-problems}
\section{Production and backreaction}
The specified rapid-crossing spectrum must be replaced or recovered by a controlled coupled field-and-mode calculation. Its energy overdraw is already a concrete obstruction, not a small unquantified uncertainty. Cutoff stability, a conserved renormalized stress tensor, and consistent propagation of the produced spectrum are prerequisites for a microscopic prediction.
\section{Microscopic derivation}
The dominant theoretical gap is the derivation of the physical refinement scale, production epoch $a_\star$, and coupling $\gamma$ from a common microscopic gravitational state.
\section{Quantum and radiative stability}
A complete theory must determine the renormalized effective action, loop stability of the hierarchy, total vacuum-energy accounting, and the fate of locking beyond the Gaussian detail approximation.
\section{Precision transport}
Required work includes precision multilevel recombination and helium kinetics, a production-grade line-of-sight solver, astrophysical reionization, massive neutrinos, full-sky lensing, and comparison with established Boltzmann codes.
\clearpage
\section{Nonlinear structure and baryons}
Required work includes larger volumes, fixed-volume resolution convergence, host zooms, six-dimensional subhalo finding, live disks, hydrodynamics, molecular and metal cooling, radiation transport, star formation, feedback, and survey-calibrated selection.
\section{Statistical establishment}
A decisive data milestone is a joint likelihood with a frozen microscopic parameter map and shared nuisance models across CMB, large-scale structure, lensing, high-redshift galaxies, reionization, and satellites.
\section{External establishment}
Internal audits and second implementations do not establish priority or empirical correctness. External mathematical review, independent numerical reproduction, a complete literature map, and observational model comparison remain separate gates.
